% Paper A — Finite Observables Do Not Determine Critical-Line Support
% Compilable LaTeX draft, generated from docs/paper-A-draft.md and the theorem files
% theorems/{B1-finite-inequality,B2-exact-collision,G-fredholm-certificate}/.
% All three theorems are Gate-A-established (independent whole-theorem review) with deposited
% independent checkers. This source matches the reviewed status exactly; the per-theorem
% statement.md files and docs/STATUS.md are authoritative.
\documentclass[11pt]{article}

\usepackage[margin=1in]{geometry}
\usepackage{amsmath,amssymb,amsthm}
\usepackage{mathtools}
\usepackage[T1]{fontenc}
\usepackage{microtype}
\usepackage{hyperref}
\hypersetup{colorlinks=true,linkcolor=blue,citecolor=blue,urlcolor=blue}

\theoremstyle{plain}
\newtheorem{theorem}{Theorem}
\newtheorem{lemma}{Lemma}
\newtheorem{proposition}{Proposition}
\newtheorem{corollary}{Corollary}
\theoremstyle{definition}
\newtheorem{definition}{Definition}
\theoremstyle{remark}
\newtheorem*{remark}{Remark}

\newcommand{\Xsym}{\mathfrak{X}_{\mathrm{sym}}}
\newcommand{\Zp}{\mathcal{Z}_+}
\newcommand{\Zm}{\mathcal{Z}_-}
\newcommand{\Ohat}{\widehat{O}}
\newcommand{\Xihat}{\widehat{\Xi}}
\newcommand{\hhat}{\widehat{h}}
\newcommand{\Md}{\mathfrak{M}}

\title{Finite Observables Do Not Determine Critical-Line Support}
\author{[authors]}
\date{Draft --- \today}

\begin{document}
\maketitle

\begin{abstract}
We prove three information-theoretic \emph{obstruction theorems} for methods that decide the
predicate ``all zeros lie on the critical line'' using only a finite or archimedean-level
observation of a zeta-like zero multiset. Let $\Xsym$ be the class of locally finite zero
multisets in the critical strip, symmetric under conjugation and $\rho\mapsto 1-\rho$, with an
admissibility exponent bound.
\begin{itemize}
  \item \textbf{Theorem~\ref{thm:B1} (no uniform separation margin).} For any finite test
  family $\Phi=(\varphi_1,\dots,\varphi_m)$ (Li, Weil--W2, or moment type) and any tolerance,
  there exist $\Zp,\Zm\in\Xsym$ with $P(\Zp)=1$, $P(\Zm)=0$, whose observation vectors satisfy
  $|O_\Phi(\Zp)-O_\Phi(\Zm)|<\varepsilon$ coordinatewise; the infimum over the construction of
  the observation gap between the $P=1$ and $P=0$ classes is $0$.
  \item \textbf{Theorem~\ref{thm:B2} (exact observation collision).} Under the same hypotheses
  one may take $O_\Phi(\Zp)=O_\Phi(\Zm)$ \emph{exactly} (zero tolerance, integer
  multiplicities), with $\Zp$ on-line and $\Zm$ additionally carrying an off-line quartet.
  \item \textbf{Theorem~\ref{thm:Ginfo} (diagonal Fredholm obstruction).} For every method in a
  natural archimedean-level class $\Md_d^{\mathrm{tr}}$, the locally uniform limit of the
  associated finite-rank Fredholm determinant is a fixed function $G_d$ that is
  unconditionally \emph{not} the target $\Xihat$.
\end{itemize}
All three exhibit a method class whose observation map does not determine the target
predicate. \textbf{These are no-go theorems for a class of methods; they do not prove,
disprove, or make progress toward the Riemann Hypothesis, and they hold regardless of its
truth.}
\end{abstract}

\section{Introduction}

\subsection{What this paper does and does not claim}
The Riemann Hypothesis (RH) is equivalent to many criteria---positivity of the Li coefficients
\cite{Li1997}, positivity of the Weil functional \cite{Weil1952,Bombieri2000}, reality of the
spectrum of conjectural operators \cite{BerryKeating1999,Connes1999}. Each such criterion
\emph{locates} the difficulty of RH in a particular object; none, by itself, is a barrier to
any proof strategy.

This paper proves the complementary kind of statement: for three precisely delimited
\emph{method classes}, the data the method observes \emph{does not determine} whether all zeros
lie on the critical line. We emphasize at the outset, and repeat in the conclusion:
\begin{quote}
\textbf{This paper does not prove, disprove, or approach RH.} Its theorems are no-go results
for method classes---statements of the form ``a method seeing only $O$ cannot decide the
predicate $P$'', proved by exhibiting two admissible objects that agree on $O$ but differ on
$P$. RH itself is never assumed, and no result here is an RH-equivalent reformulation.
\end{quote}
We call such a theorem a \emph{barrier} for a method class only when all five of the following
are explicit: the method class, the ambient object class, the observation map, the target
predicate, and an escape route (a strictly larger class or richer observation to which the
obstruction does not apply). We are careful to distinguish a genuine obstruction from the
non-results the program excludes: another RH-equivalence; a finite failure of one sufficient
inequality; a positive margin tending to zero; or a synthetic configuration outside the
ambient class.

\subsection{Relation to prior work}
That finite real-zero data cannot certify an infinite positivity criterion is folklore, and
finite-Li limitations were noted by Bombieri--Lagarias \cite{BombieriLagarias1999}. Our
contribution is (i) to make the obstruction \emph{exact}---zero tolerance and integer
multiplicities (Theorem~\ref{thm:B2})---in a precisely defined ambient class with an explicit
escape-route list; (ii) to prove the companion \emph{no uniform margin} statement
(Theorem~\ref{thm:B1}) that keeps the result clear of the ``margin $\to 0$'' non-barrier; and
(iii) to give the analogous archimedean-level obstruction (Theorem~\ref{thm:Ginfo}) for the
diagonal Fredholm-determinant approach in the sense of Connes--Consani--Moscovici
\cite{CCM2025}.

\section{Setup}\label{sec:setup}

\begin{definition}[Ambient class]
$\Xsym$ is the set of locally finite multisets $\mathcal{Z}$ of points in the open critical
strip $\{0<\operatorname{Re} s<1\}$, symmetric under $\rho\mapsto\bar\rho$ and $\rho\mapsto 1-\rho$, with a
fixed admissibility exponent: $\sum_{\rho\in\mathcal Z}|\rho|^{-(1+\varepsilon)}<\infty$ for
some $\varepsilon>0$. Membership is checkable, and any finite symmetric multiset is a member.
\end{definition}

\begin{definition}[Target predicate]
$P(\mathcal Z)=1$ iff every atom $\rho\in\mathcal Z$ has $\operatorname{Re}\rho=\tfrac12$; else
$P(\mathcal Z)=0$.
\end{definition}

\paragraph{Test families.} We use three interchangeable families:
Li-type $\varphi_j(\rho)=1-(1-1/\rho)^j$; Weil--W2 $\varphi_j(\rho)=\hhat_j(\operatorname{Im}\rho)$ with
$h_j\in C_c^\infty(\mathbb R;\mathbb R)$ even and
$\hhat_j(\xi)=\int h_j(x)e^{i\xi x}\,dx$ (evaluation at the imaginary part only---complex-point
Weil evaluation is excluded, as $\hhat$ has exponential type and need not decay); and moment
$\varphi_j(\rho)=\rho^{-j}$, $j\ge 1$.

\paragraph{Observation maps and the $\Sigma'$ convention.} The two finite-observation theorems
use \emph{different} summation conventions, differing by a factor of $2$; a numerical anchor
must not be transported between them.
\begin{itemize}
  \item \textbf{B1 (R-atom):} $O_j(\mathcal Z)=\sum'_{\rho\in\mathcal Z}\varphi_j(\rho)$---each
  atom once, the prime a symmetric-height regularization for convergence, \emph{not} doubling.
  \item \textbf{B2 (R-symm):} $O_j(\mathcal Z)=\sum_{\rho\in\mathcal
  Z}[\varphi_j(\rho)+\varphi_j(1-\rho)]$---which doubles on a multiset closed under
  $\rho\mapsto 1-\rho$.
\end{itemize}
Both give $O_j\in\mathbb R$ for the real-coefficient families. B2's collision identity is
scale-invariant (\S\ref{sec:B2}), so the factor is harmless there; B1's decay anchors are
stated in the R-atom convention and are \emph{not} the R-symm values.

\paragraph{Off-line quartet.} For $\sigma_0\in(0,1)\setminus\{\tfrac12\}$ and $T>0$,
$Q(\sigma_0,T):=\{\sigma_0+iT,\,1-\sigma_0+iT,\,\sigma_0-iT,\,1-\sigma_0-iT\}$, a symmetric
member of $\Xsym$ with $P(Q)=0$. We fix $\sigma_0=\tfrac34$ throughout.

\section{Theorem B1 --- no uniform separation margin}\label{sec:B1}

\begin{theorem}\label{thm:B1}
Let $\Phi=(\varphi_1,\dots,\varphi_m)$ be any finite test family of the above types, let
$\Zp\in\Xsym$ with $P(\Zp)=1$ satisfy strict inequalities $(O_\Phi(\Zp))_j>c_j$
($j=1,\dots,m$), and let $\varepsilon_j>0$. Then there exists $\Zm\in\Xsym$ with $P(\Zm)=0$ and
$|(O_\Phi(\Zm))_j-(O_\Phi(\Zp))_j|<\varepsilon_j$ for all $j$; in particular
$(O_\Phi(\Zm))_j>c_j$ whenever $\varepsilon_j<(O_\Phi(\Zp))_j-c_j$. Consequently the infimum,
over admissible constructions, of the coordinatewise observation gap between the $P=1$ and
$P=0$ classes is $0$: no fixed positive coordinate margin separates the two classes.
\end{theorem}

\noindent($\Zp$ ranges over abstract $P=1$ members---e.g.\ explicit finite on-line multisets.
The instantiation ``$\Zp=$ the $\zeta$ zero multiset'' is conditional on RH and is not used.)

\begin{lemma}[Off-line quartet]\label{lem:B1}
For $\varphi$ of any of the three types and $\sigma_0\in(0,1)\setminus\{\tfrac12\}$, the
quartet contribution
$\delta_j(T):=\varphi_j(\sigma_0+iT)+\varphi_j(1-\sigma_0+iT)+\varphi_j(\sigma_0-iT)+\varphi_j(1-\sigma_0-iT)$
is continuous in $T>0$ and $\delta_j(T)\to 0$ as $T\to\infty$.
\end{lemma}
\begin{proof}
\emph{Li-type:} $|1/\rho|\le(T^2+\sigma_0^2)^{-1/2}\to 0$, so $\varphi_j(\sigma_0\pm iT)\to 0$
and likewise for $1-\sigma_0\pm iT$. \emph{Weil--W2:} $h_j\in C_c^\infty$ gives, by $N$-fold
integration by parts, $|\hhat_j(T)|\le C_N(1+|T|)^{-N}$; with $h_j$ even,
$\delta_j(T)=4\hhat_j(T)\to 0$ (Riemann--Lebesgue, \cite[Thm.~9.6]{Rudin1987}).
\emph{Moment:} $|\varphi_j(\rho)|=|\rho|^{-k}=O(T^{-k})\to 0$. Continuity is immediate from
continuity of $\varphi_j$ in $\rho$ and of the quartet points in $T$.
\end{proof}

\begin{remark}[Convergence of $\Sigma'$, R-atom]
The Li terms are only $O(|\rho|^{-1})$, not dominated by the admissibility series
$\sum|\rho|^{-(1+\eta)}$. Grouping each $\rho=\sigma+it$ with its conjugate, the
real-coefficient pair contributes $2\operatorname{Re}\varphi_j$; the leading term gives
$2\operatorname{Re}(1/\rho)=2\sigma/(\sigma^2+t^2)=O(|\operatorname{Im}\rho|^{-2})$, and every higher inverse power is also
$O(|\operatorname{Im}\rho|^{-2})$; since $1+\eta<2$, the grouped sum converges absolutely.
\end{remark}

\begin{proof}[Proof of Theorem~\ref{thm:B1}]
By Lemma~\ref{lem:B1} choose $T_*$ so large that $|\delta_j(T_*)|<\varepsilon_j$ for all
$j=1,\dots,m$ (possible since $m$ is finite). Set $\Zm:=\Zp\cup Q(\tfrac34,T_*)$. Then
$\Zm\in\Xsym$ (union of admissible members); $P(\Zm)=0$ (the atoms $\tfrac34\pm iT_*$ are
off-line); and $|(O_\Phi(\Zm))_j-(O_\Phi(\Zp))_j|=|\delta_j(T_*)|<\varepsilon_j$.
\end{proof}

\subsection{Precise strength}
Theorem~\ref{thm:B1} asserts that the observation-gap \emph{infimum} between the classes is
$0$---there is no positive uniform margin. It does \emph{not} assert an exact-discriminator
failure at fixed inputs; distinguishing the two classes by an exact discontinuous rule on
precise real inputs is not excluded by B1 (that is B2's job). The quartet contribution is
strictly positive for each finite $T$ (e.g.\
$\delta_1(T)=2[\sigma/(\sigma^2+T^2)+(1-\sigma)/((1-\sigma)^2+T^2)]>0$ for the Li first
coordinate), so this is a statement about the observation-separation infimum, \emph{not} a
shrinking positivity margin of a sufficient inequality---keeping B1 clear of the
``margin $\to 0$'' non-barrier. In the R-atom convention the exact decay anchors are
$\delta_1(1)=\tfrac{608}{425}$, threshold $T_*=90$ (for $\sigma_0=\tfrac34$, $m=2$,
$\varepsilon=10^{-3}$), and $\delta_j(T)\cdot T^2\to 2j^2$; these are independently
machine-checked (Appendix~\ref{app:checkers}).

\section{Theorem B2 --- exact observation collision}\label{sec:B2}

\begin{theorem}\label{thm:B2}
For any finite test family $\Phi=(\varphi_1,\dots,\varphi_m)$ of Li or moment type there exist
$\Zp,\Zm\in\Xsym$ with $P(\Zp)=1$, $P(\Zm)=0$, all multiplicities integers, and
$O_\Phi(\Zp)=O_\Phi(\Zm)$ exactly. Here $\Zp$ consists of on-line pairs and $\Zm$ differs by
adding an off-line quartet and adjusting on-line multiplicities.
\end{theorem}

We use the R-symm convention, additive under multiset union; on one on-line pair
$L(t)=\{\tfrac12\pm it\}$ it gives $O_j(L(t))=4\operatorname{Re}\varphi_j(\tfrac12+it)$, and on the quartet
$O_j(Q(\tfrac34,T))=4\operatorname{Re}[\varphi_j(\tfrac34+iT)+\varphi_j(\tfrac14+iT)]$.

\subsection{Nonsingular Jacobian (self-contained Vandermonde)}
Fix distinct rational heights $t_1<\dots<t_m>0$. For the Li family, with
$x_k=(4t_k^2-1)/(4t_k^2+1)\in\mathbb Q$ and Chebyshev $T_j$,
\[
  C_{jk}:=O_j(L(t_k))=4\bigl(1-T_j(x_k)\bigr),\qquad j,k=1,\dots,m.
\]
Writing $1-T_j(x)=(1-x)Q_{j-1}(x)$ with $\deg Q_{j-1}=j-1$ and leading coefficient $2^{j-1}$,
the matrix $[Q_{j-1}(x_k)]$ equals $U\cdot V(x_1,\dots,x_m)$ for an upper-triangular $U$
(Chebyshev-to-monomial change of basis) and the Vandermonde $V$. Hence
\[
  \det C = 4^m\Bigl(\textstyle\prod_k(1-x_k)\Bigr)\,2^{m(m-1)/2}\!\!\prod_{k<l}(x_l-x_k)\neq 0,
\]
all factors nonzero ($1-x_k>0$; the $x_k$ distinct and increasing). For rational $t_k$,
$x_k\in\mathbb Q$, so $C\in\mathbb Q^{m\times m}$ and $\det C\in\mathbb Q\setminus\{0\}$. (The
moment family gives the analogous cosine-Vandermonde $[\cos(j\phi_k)]$, nonsingular by the same
reduction.)

\subsection{Exact collision with integer multiplicities}
Let $d(T)_j:=O_j(Q(\tfrac34,T))\in\mathbb Q^m$ (rational for rational $T$) and solve
$C\alpha^{\mathbb Q}=-d(T)$, $\alpha^{\mathbb Q}\in\mathbb Q^m$. Put
$R:=\operatorname{lcm}(\text{denominators of }\alpha^{\mathbb Q})$ and
$n_k:=R\,\alpha^{\mathbb Q}_k\in\mathbb Z$. Scaling by $R$ gives the exact integer identity
\[
  C\,n + R\,d(T)=0.
\]
This identity is scale-invariant ($C\to 2C$, $d\to 2d$ leaves $n,R$ and the identity
unchanged), so the R-symm-vs-R-atom factor is immaterial here. Set the multiplicity buffer
$M:=\max_k|n_k|$ and define
\[
  \Zp:=\bigsqcup_{k=1}^m M\cdot L(t_k),\qquad
  \Zm:=\bigsqcup_{k=1}^m (M+n_k)\cdot L(t_k)\;\sqcup\; R\cdot Q(\tfrac34,T).
\]
Both are in $\Xsym$; $M+n_k\ge 0$ so $\Zm$ is a valid multiset; $P(\Zp)=1$, $P(\Zm)=0$. Since
$O_\Phi$ is additive,
\[
  O_j(\Zm)-O_j(\Zp)=\sum_k n_k\,O_j(L(t_k))+R\,O_j(Q(\tfrac34,T))=(C\,n)_j+R\,d(T)_j=0. \qedhere
\]

\subsection{Non-vacuity and scope}
Both adversaries are \emph{constructed} finite members of $\Xsym$; B2 does \emph{not} assert
that $\zeta$'s zeros are indistinguishable from an off-line configuration. A structural
non-vacuity check---the quartet's first coordinate
$d_1(T)=64(16T^2+3)/(256T^4+160T^2+9)>0$ for all real $T$---confirms the target vector is
nonzero, so the solve is nontrivial. The exact-rational collision has been independently
reconstructed from the definition (two agreeing routes, adversarial mutation guard);
Appendix~\ref{app:checkers}.

\section{Theorem G-info --- diagonal Fredholm obstruction}\label{sec:Ginfo}

We pass one observation layer up, to the archimedean levels used in
Connes--Consani--Moscovici-style spectral-determinant approaches \cite{CCM2025}.

\paragraph{Observation.} $O_\theta(n)=d_n:=\theta_{\mathrm{level}}(n)$, the zero-free
Riemann--Siegel unfolding levels
($\theta(t)=\operatorname{Im}\log\Gamma(\tfrac14+\tfrac{it}{2})-\tfrac{t}{2}\log\pi$), determined by the
$\Gamma$-function alone.
\paragraph{Method class.} $\Md_d^{\mathrm{tr}}$: finite-rank positive semidefinite families
$(K_N)$ with $K_N=\Phi_N(d_1,\dots,d_N)$ and $\|K_N-D_N\|_1\to 0$, where
$D_N=\operatorname{diag}(\kappa_1,\dots,\kappa_N)$, $\kappa_n=1/(\tfrac14+d_n^2)$. Non-empty
(witness $K_N=D_N$).
\paragraph{Target.} The predicate is ``the locally uniform limit of $\det(I-z^2K_N)$ equals
$\Xihat$'', the entire CCM target (normalized $\Xihat(0)=1$), whose zeros are real iff RH holds.

\begin{theorem}\label{thm:Ginfo}
For every $(K_N)\in\Md_d^{\mathrm{tr}}$,
\[
  \det(I-z^2K_N)\;\longrightarrow\;G_d(z)=\prod_{n\ge 1}\Bigl(1-\frac{z^2}{\tfrac14+d_n^2}\Bigr)
\]
locally uniformly, and $G_d\neq\Xihat$ unconditionally.
\end{theorem}
\begin{proof}
Locally uniform convergence follows from trace-norm stability of Fredholm determinants and
$\|K_N-D_N\|_1\to 0$. For $G_d\neq\Xihat$: if RH holds, $\Xihat=C\prod(1-z^2/\gamma_n^2)$ has
zeros $\{\pm\gamma_n\}$, while $G_d$ has zeros $\{\pm\sqrt{\tfrac14+d_n^2}\}$; these differ
because $\sqrt{\tfrac14+d_n^2}>d_n$ and $d_n\neq\gamma_n$ for infinitely many $n$ (the
arithmetic-fluctuation gap $\gamma_n-d_n\sim S(\gamma_n)/A'(\gamma_n)$, $A'(t)=\theta'(t)/\pi$;
proved unconditionally). If RH fails, $\Xihat$ has a non-real zero while every zero of $G_d$ is
real. Either way $G_d\neq\Xihat$.
\end{proof}

Two independent separations are visible: (i) a spectral-parameter shift
$d_n\mapsto\sqrt{\tfrac14+d_n^2}$ (the diagonal determinant's zeros are not at $\pm d_n$), and
(ii) the $S(T)$ gap $\{d_n\}\neq\{\gamma_n\}$; the first survives even if one corrects the
$\tfrac14$ shift.

\paragraph{Scope: G-info vs.\ G-hard.} Theorem~\ref{thm:Ginfo} is the diagonal obstruction.
The general claim ``no method in the larger class $\Md_{\mathrm{FC}}$ can supply the $S(T)$
data without reading zero ordinates or computing an RH-equivalent quantity'' is a separate
\emph{conjecture} (G-hard) and is \emph{not} claimed here.

\paragraph{Why convergence would be RH-strength.} If some $(K_N)$ in the class had
$\det(I-z^2K_N)\to\Xihat$ locally uniformly, then by Hurwitz the real zeros of the approximants
would force $\Xihat$ to have all real zeros---i.e.\ RH. So a convergence claim in this class
cannot be a zero-independent input; it would \emph{imply} RH. Theorem~\ref{thm:Ginfo} shows the
diagonal limit is a different function, $G_d$, so the approach does not even converge to the
target.

\section{Escape routes}\label{sec:escape}
Each theorem names a strictly larger class or richer observation to which it does \emph{not}
apply:
\begin{enumerate}
  \item Infinite test hierarchy ($K\to\infty$, all Li/Weil data).
  \item Euler product / multiplicative structure (Selberg-class axioms).
  \item Gamma factor and analytic continuation---archimedean data beyond $O_\Phi$.
  \item Coefficient arithmetic (integrality/positivity of Dirichlet coefficients).
  \item For G-info: methods supplying genuine $S(T)$ data (the G-hard frontier), or
  non-diagonal $K_N$ outside $\Md_d^{\mathrm{tr}}$.
\end{enumerate}

\section{Limitations and honest boundaries}\label{sec:limits}
\begin{itemize}
  \item \textbf{Not RH.} No theorem here proves, disproves, or approaches RH; RH is neither
  assumed nor concluded. B1/B2 construct adversary multisets, not $\zeta$; G-info's
  $\mathcal Z_{\mathrm{RH}}$ appears only as a hypothetical target and its reality is never
  used.
  \item B1 is ``no uniform margin'', not an exact-discriminator failure (\S\ref{sec:B1}).
  \item B2's adversaries are constructed, not $\zeta$'s zeros (\S\ref{sec:B2}).
  \item G-info is the diagonal obstruction only; G-hard is a conjecture (\S\ref{sec:Ginfo}).
  \item Finite certificates validate finite statements, not the analytic theorems.
\end{itemize}

\appendix
\section{Deposited independent checkers}\label{app:checkers}
All are pure-standard-library, exact-rational or outward-rounded interval, with no floating
point in the certificate path; each re-runs to \texttt{ALL\_CERTIFIED\_CHECKS\_PASSED} and is
pinned by SHA-256.

\begin{center}
\begin{tabular}{lll}
\hline
Checker & Theorem & SHA-256 (prefix) \\
\hline
\texttt{b1\_ratom\_certified\_checker.py} & B1 & \texttt{199c7dad}\\
\texttt{b2\_certified\_checker.py} & B2 & \texttt{776eeab5}\\
\texttt{diagonal\_fredholm\_interval\_replay.py} & G-info & \texttt{e197f2bb}\\
\hline
\end{tabular}
\end{center}
\noindent They validate, respectively: the R-atom decay
($\delta_1(1)=\tfrac{608}{425}$, $T_*=90$, $\delta_jT^2\to 2j^2$); the exact collision
$C\,n+R\,d=0$ via two agreeing routes with an adversarial mutation guard; and the Gram levels
$d_n$ with the separation $\gamma_n<d_n<\sqrt{\tfrac14+d_n^2}$ and the tail bound. A finite
certificate validates only the finite replayed statement, never the analytic theorem.

\begin{thebibliography}{9}
\bibitem{BerryKeating1999} M.~V. Berry and J.~P. Keating, \emph{$H=xp$ and the Riemann zeros},
  SIAM Review \textbf{41} (1999), 236--266.
\bibitem{Bombieri2000} E.~Bombieri, \emph{The Riemann Hypothesis}, Clay Mathematics Institute
  problem description, 2000.
\bibitem{BombieriLagarias1999} E.~Bombieri and J.~C. Lagarias, \emph{Complements to Li's
  criterion for the Riemann Hypothesis}, J.~Number Theory \textbf{77} (1999), 274--287.
\bibitem{CCM2025} A.~Connes, C.~Consani, and H.~Moscovici, \emph{Zeta spectral triples},
  arXiv:2511.22755 (2025).
\bibitem{Connes1999} A.~Connes, \emph{Trace formula in noncommutative geometry and the zeros of
  the Riemann zeta function}, Selecta Math. \textbf{5} (1999), 29--106.
\bibitem{Li1997} X.-J. Li, \emph{The positivity of a sequence of numbers and the Riemann
  Hypothesis}, J.~Number Theory \textbf{65} (1997), 325--333.
\bibitem{Rudin1987} W.~Rudin, \emph{Real and Complex Analysis}, 3rd ed., McGraw-Hill, 1987.
\bibitem{Weil1952} A.~Weil, \emph{Sur les ``formules explicites'' de la th\'eorie des nombres
  premiers}, Comm. S\'em. Math. Univ. Lund (1952), 252--265.
\end{thebibliography}

\end{document}
